% 无穷多个阈值在四个训练分数上只给出五种预测模式（第5章5.2.1）。
% 训练分数 0.1、0.3、0.7、0.9；阈值规则 h_t(x)=I{x>=t}，t in [0,1]。
% 小数轴映射：屏幕横坐标 = 34 + 50*s（s 为分数），每行用色带标出给出该模式的阈值段。
% 实心蓝点为预测 1，空心红圈为预测 0；第三行另标 t=0.5 与 t=0.6 两个具体阈值。
\begingroup
\providecolor{FigureGreen}{HTML}{829C8B}
\providecolor{FigureGreenFill}{HTML}{E8EEE9}
\begin{tikzpicture}[figure picture,x=1mm,y=1mm,
  one/.style={circle,fill=FigureBlue,draw=FigureBlue,inner sep=0pt,minimum size=2.6mm},
  zero/.style={circle,fill=white,draw=FigureRed,line width=.75pt,inner sep=0pt,minimum size=2.6mm}]
\node[figure heading,anchor=west] at (0,49) {阈值所在区间};
\node[figure heading] at (59,49) {四封邮件上的预测};
\node[figure heading] at (98,49) {模式};
\foreach \x/\s in {39/0.1,49/0.3,69/0.7,79/0.9}{\node[figure note] at (\x,44) {$\s$};}
\foreach \row/\lab/\a/\b/\pa/\pb/\pc/\pd/\pat in {
  0/{$t\leq0.1$}/34/39/one/one/one/one/{(1,1,1,1)},
  1/{$0.1<t\leq0.3$}/39/49/zero/one/one/one/{(0,1,1,1)},
  2/{$0.3<t\leq0.7$}/49/69/zero/zero/one/one/{(0,0,1,1)},
  3/{$0.7<t\leq0.9$}/69/79/zero/zero/zero/one/{(0,0,0,1)},
  4/{$0.9<t\leq1$}/79/84/zero/zero/zero/zero/{(0,0,0,0)}}{
  \pgfmathsetmacro{\y}{38-9*\row}
  \node[anchor=west] at (0,\y) {\lab};
  \fill[FigureGreenFill] (\a,\y-2.2) rectangle (\b,\y+2.2);
  \draw[FigureGreen,line width=.9pt] (\a,\y+2.2)--(\b,\y+2.2);
  \draw[FigureStroke] (34,\y)--(84,\y);
  \node[\pa] at (39,\y) {};\node[\pb] at (49,\y) {};
  \node[\pc] at (69,\y) {};\node[\pd] at (79,\y) {};
  \node at (98,\y) {\pat};
}
\foreach \x/\t in {59/{0.5},64/{0.6}}{
  \draw[FigureInk,line width=.75pt] (\x,18.6)--(\x,21.4);
  \node[figure note] at (\x,16) {$\t$};}
\end{tikzpicture}
\endgroup
